% Môn: Vật lý
% Lớp: 12
% Chương: Chương II. Khí lí tưởng
% Bài: L12C2 Bài 8. Mô hình động học phân tử chất khí · Khối khí chịu áp suất của cột chất lỏng
% ===== Câu 18 =====
% ID: L12C2B8-93-DS
% Mức: TH
\dangbt{Khối khí chịu áp suất của cột chất lỏng}
\begin{ex}
% ID: L12C2B8-93-DS
% Mức: TH
Trong dạng “Khối khí chịu áp suất của cột chất lỏng”, Trình bày cơ sở lí thuyết của dạng “Khối khí chịu áp suất của cột chất lỏng” và nêu điều kiện áp dụng.
\choiceTF

\loigiai{
Nêu đúng đại lượng không đổi, hệ thức đặc trưng và yêu cầu dùng nhiệt độ kelvin, áp suất tuyệt đối khi có liên quan.
}
\end{ex}

% ===== Câu 22 =====
% ID: L12C2B8-94-DS
% Mức: TH
\begin{ex}
% ID: L12C2B8-94-DS
% Mức: TH
\choiceTF{\True Trong quá trình đẳng nhiệt, đồ thị biểu diễn mối quan hệ giữa áp suất $P$ và thể tích $V$ là một đường hypebol.}{\True Trong quá trình đẳng áp, đồ thị biểu diễn mối quan hệ giữa thể tích $V$ và nhiệt độ tuyệt đối $T$ là một đường thẳng đi qua gốc tọa độ.}{Trong quá trình đẳng tích, đồ thị biểu diễn mối quan hệ giữa áp suất $P$ và nhiệt độ tuyệt đối $T$ là một đường thẳng không đi qua gốc tọa độ.}{\True Để nhận biết quá trình đẳng nhiệt, cần xác định nhiệt độ của khối khí không đổi trong suốt quá trình.}


\loigiai{
a) Phát biểu này đúng. Quá trình đẳng nhiệt tuân theo định luật Boyle-Mariotte: $PV = \text{const}$. Đồ thị của hàm số này trên hệ trục tọa độ $P-V$ là một đường hypebol.
b) Phát biểu này đúng. Quá trình đẳng áp tuân theo định luật Gay-Lussac: $\frac{V}{T} = \text{const}$. Đồ thị của hàm số này trên hệ trục tọa độ $V-T$ là một đường thẳng đi qua gốc tọa độ (nếu $V$ và $T$ là các giá trị tuyệt đối).
c) Phát biểu này sai. Quá trình đẳng tích tuân theo định luật Charles: $\frac{P}{T} = \text{const}$. Đồ thị của hàm số này trên hệ trục tọa độ $P-T$ là một đường thẳng đi qua gốc tọa độ (nếu $P$ và $T$ là các giá trị tuyệt đối).
d) Phát biểu này đúng. Để nhận biết quá trình đẳng nhiệt, điều kiện tiên quyết là nhiệt độ của khối khí phải không đổi. Điều này có thể được kiểm tra bằng cách đo nhiệt độ hoặc quan sát các yếu tố môi trường ảnh hưởng đến nhiệt độ.
}
\end{ex}

% ===== Câu 29 =====
% ID: L12C2B8-95-DS
% Mức: TH
\begin{ex}
% ID: L12C2B8-95-DS
% Mức: TH
Xét các phát biểu sau trong dạng “Khối khí chịu áp suất của cột chất lỏng”.
\choiceTF
{\True Đại lượng nào không đổi trong quá trình đẳng nhiệt của một lượng khí xác định? — phát biểu đúng là: Nhiệt độ.}
{Đại lượng nào không đổi trong quá trình đẳng nhiệt của một lượng khí xác định? — phát biểu Áp suất là đúng.}
{\True Hệ thức nào là định luật Boyle? — phát biểu đúng là: $p_1V_1=p_2V_2$.}
{Hệ thức nào là định luật Boyle? — phát biểu $V_1/T_1=V_2/T_2$ là đúng.}
\loigiai{
\begin{itemchoice}
\item (Đúng) Đại lượng nào không đổi trong quá trình đẳng nhiệt của một lượng khí xác định? — phát biểu đúng là: Nhiệt độ. (Vì): Quá trình đẳng nhiệt có nhiệt độ không đổi. b) Sai: trái với hệ thức hoặc điều kiện đã nêu. c) Đúng: Ở nhiệt độ không đổi, $pV$ là hằng số. d) Sai: trái với hệ thức hoặc điều kiện đã nêu
\item (Sai) Đại lượng nào không đổi trong quá trình đẳng nhiệt của một lượng khí xác định? — phát biểu Áp suất là đúng.
\item (Đúng) Hệ thức nào là định luật Boyle? — phát biểu đúng là: $p_1V_1=p_2V_2$.
\item (Sai) Hệ thức nào là định luật Boyle? — phát biểu $V_1/T_1=V_2/T_2$ là đúng.
\end{itemchoice}
}
\end{ex}

% ===== Câu 33 =====
% ID: L12C2B8-96-DS
% Mức: TH
\begin{ex}
% ID: L12C2B8-96-DS
% Mức: TH
Xét các phát biểu sau trong dạng “Khối khí chịu áp suất của cột chất lỏng”.
\choiceTF
{\True Khí có áp suất $120\,\text{kPa}$, thể tích 12 lít được nén đẳng nhiệt còn 3 lít. Áp suất sau là — phát biểu đúng là: $480\,\text{kPa}$.}
{Khí có áp suất $120\,\text{kPa}$, thể tích 12 lít được nén đẳng nhiệt còn 3 lít. Áp suất sau là — phát biểu $30\,\text{kPa}$ là đúng.}
{\True Khí có áp suất $125\,\text{kPa}$, thể tích 5 lít giãn đẳng nhiệt đến 25 lít. Áp suất sau là — phát biểu đúng là: $25\,\text{kPa}$.}
{Khí có áp suất $125\,\text{kPa}$, thể tích 5 lít giãn đẳng nhiệt đến 25 lít. Áp suất sau là — phát biểu $125\,\text{kPa}$ là đúng.}
\loigiai{
\begin{itemchoice}
\item (Đúng) Khí có áp suất $120\,\text{kPa}$, thể tích 12 lít được nén đẳng nhiệt còn 3 lít. Áp suất sau là — phát biểu đúng là: $480\,\text{kPa}$. (Vì): $p_2=p_1V_1/V_2=120\cdot12/3=480$ kPa. b) Sai: trái với hệ thức hoặc điều kiện đã nêu. c) Đúng: $p_2=125\cdot5/25=25$ kPa. d) Sai: trái với hệ thức hoặc điều kiện đã nêu
\item (Sai) Khí có áp suất $120\,\text{kPa}$, thể tích 12 lít được nén đẳng nhiệt còn 3 lít. Áp suất sau là — phát biểu $30\,\text{kPa}$ là đúng.
\item (Đúng) Khí có áp suất $125\,\text{kPa}$, thể tích 5 lít giãn đẳng nhiệt đến 25 lít. Áp suất sau là — phát biểu đúng là: $25\,\text{kPa}$.
\item (Sai) Khí có áp suất $125\,\text{kPa}$, thể tích 5 lít giãn đẳng nhiệt đến 25 lít. Áp suất sau là — phát biểu $125\,\text{kPa}$ là đúng.
\end{itemchoice}
}
\end{ex}

% ===== Câu 37 =====
% ID: L12C2B8-97-DS
% Mức: VD
\begin{ex}
% ID: L12C2B8-97-DS
% Mức: VD
Xét các phát biểu sau trong dạng “Khối khí chịu áp suất của cột chất lỏng”.
\choiceTF
{\True Thể tích khí giảm 2 lần trong quá trình đẳng nhiệt thì áp suất — phát biểu đúng là: tăng 2 lần.}
{Thể tích khí giảm 2 lần trong quá trình đẳng nhiệt thì áp suất — phát biểu giảm 2 lần là đúng.}
{\True Trên hệ trục p- $V$, đường đẳng nhiệt là — phát biểu đúng là: một nhánh hypebol.}
{Trên hệ trục p- $V$, đường đẳng nhiệt là — phát biểu đường thẳng song song trục p là đúng.}
\loigiai{
\begin{itemchoice}
\item (Đúng) Thể tích khí giảm 2 lần trong quá trình đẳng nhiệt thì áp suất — phát biểu đúng là: tăng 2 lần. (Vì): o $p$ tỉ lệ nghịch với $V$. b) Sai: trái với hệ thức hoặc điều kiện đã nêu. c) Đúng: Từ $p=C/V$, đồ thị là một nhánh hypebol. d) Sai: trái với hệ thức hoặc điều kiện đã nêu
\item (Sai) Thể tích khí giảm 2 lần trong quá trình đẳng nhiệt thì áp suất — phát biểu giảm 2 lần là đúng.
\item (Đúng) Trên hệ trục p- $V$, đường đẳng nhiệt là — phát biểu đúng là: một nhánh hypebol.
\item (Sai) Trên hệ trục p- $V$, đường đẳng nhiệt là — phát biểu đường thẳng song song trục p là đúng.
\end{itemchoice}
}
\end{ex}

% ===== Câu 38 =====
% ID: L12C2B8-98-TL
% Mức: VDC
\begin{bt}
% ID: L12C2B8-98-TL
% Mức: VDC
Trong dạng “Khối khí chịu áp suất của cột chất lỏng”, undefined
\loigiai{
$V_2=130\cdot10/(260)=5$ lít.
}
\end{bt}

% ===== Câu 40 =====
% ID: L12C2B8-99-TN
% Mức: VDC
\begin{ex}
% ID: L12C2B8-99-TN
% Mức: VDC
Trong dạng “Bài toán tổng hợp về quá trình đẳng nhiệt”, Trong ống có cột chất lỏng, áp suất khí bị giam được xác định bằng
\choice
{\True cân bằng áp suất tại cùng độ cao}
{chỉ áp suất khí quyển}
{chỉ $\rho gh$}
{luôn bằng không}
\loigiai{
Đáp án A. Dùng $p_{khí}=p_0\pm\rho gh$ tùy cách bố trí.
}
\end{ex}

% ===== Câu 41 =====
% ID: L12C2B8-100-DS
% Mức: VD
\begin{ex}
% ID: L12C2B8-100-DS
% Mức: VD
Xét các phát biểu sau trong dạng “Khối khí chịu áp suất của cột chất lỏng”.
\choiceTF
{\True Khi dùng định luật Boyle phải sử dụng — phát biểu đúng là: áp suất tuyệt đối.}
{Khi dùng định luật Boyle phải sử dụng — phát biểu áp suất dư là đúng.}
{\True Muốn nén khí gần đẳng nhiệt nên thực hiện — phát biểu đúng là: chậm để khí trao đổi nhiệt.}
{Muốn nén khí gần đẳng nhiệt nên thực hiện — phát biểu ở thể tích không đổi là đúng.}
\loigiai{
\begin{itemchoice}
\item (Đúng) Khi dùng định luật Boyle phải sử dụng — phát biểu đúng là: áp suất tuyệt đối. (Vì): Phương trình trạng thái dùng áp suất tuyệt đối của khí. b) Sai: trái với hệ thức hoặc điều kiện đã nêu. c) Đúng: Nén chậm giúp nhiệt độ khí gần như không đổi. d) Sai: trái với hệ thức hoặc điều kiện đã nêu
\item (Sai) Khi dùng định luật Boyle phải sử dụng — phát biểu áp suất dư là đúng.
\item (Đúng) Muốn nén khí gần đẳng nhiệt nên thực hiện — phát biểu đúng là: chậm để khí trao đổi nhiệt.
\item (Sai) Muốn nén khí gần đẳng nhiệt nên thực hiện — phát biểu ở thể tích không đổi là đúng.
\end{itemchoice}
}
\end{ex}

% ===== Câu 45 =====
% ID: L12C2B8-101-DS
% Mức: VDC
\begin{ex}
% ID: L12C2B8-101-DS
% Mức: VDC
Xét các phát biểu sau trong dạng “Khối khí chịu áp suất của cột chất lỏng”.
\choiceTF
{\True Với lượng khí cố định đẳng nhiệt, khối lượng riêng tỉ lệ — phát biểu đúng là: thuận với áp suất.}
{Với lượng khí cố định đẳng nhiệt, khối lượng riêng tỉ lệ — phát biểu nghịch với áp suất là đúng.}
{\True Trong ống có cột chất lỏng, áp suất khí bị giam được xác định bằng — phát biểu đúng là: cân bằng áp suất tại cùng độ cao.}
{Trong ống có cột chất lỏng, áp suất khí bị giam được xác định bằng — phát biểu chỉ $\rho gh$ là đúng.}
\loigiai{
\begin{itemchoice}
\item (Đúng) Với lượng khí cố định đẳng nhiệt, khối lượng riêng tỉ lệ — phát biểu đúng là: thuận với áp suất. (Vì): Vì $\rho=m/V$ và $pV$ không đổi nên $\rho\sim p$. b) Sai: trái với hệ thức hoặc điều kiện đã nêu. c) Đúng: Dùng $p_{khí}=p_0\pm\rho gh$ tùy cách bố trí. d) Sai: trái với hệ thức hoặc điều kiện đã nêu
\item (Sai) Với lượng khí cố định đẳng nhiệt, khối lượng riêng tỉ lệ — phát biểu nghịch với áp suất là đúng.
\item (Đúng) Trong ống có cột chất lỏng, áp suất khí bị giam được xác định bằng — phát biểu đúng là: cân bằng áp suất tại cùng độ cao.
\item (Sai) Trong ống có cột chất lỏng, áp suất khí bị giam được xác định bằng — phát biểu chỉ $\rho gh$ là đúng.
\end{itemchoice}
}
\end{ex}

% ===== Câu 58 =====
% ID: L12C2B8-102-DS
% Mức: VD
\begin{ex}
% ID: L12C2B8-102-DS
% Mức: VD
Trong dạng “Nhận biết quá trình đẳng nhiệt và định luật Boyle”, Nêu cách xác định áp suất cột khí bị giam bởi chất lỏng.
\choiceTF

\loigiai{
Lập cân bằng áp suất tại cùng một độ cao trong chất lỏng; tùy vị trí mặt thoáng, $p_{khí}=p_0\pm\rho gh$.
}
\end{ex}

% ===== Câu 82 =====
% ID: L12C2B8-103-DS
% Mức: VD
\begin{ex}
% ID: L12C2B8-103-DS
% Mức: VD
Trong dạng “Tính áp suất, thể tích trong quá trình đẳng nhiệt”, Nêu cách xác định áp suất cột khí bị giam bởi chất lỏng.
\choiceTF

\loigiai{
Lập cân bằng áp suất tại cùng một độ cao trong chất lỏng; tùy vị trí mặt thoáng, $p_{khí}=p_0\pm\rho gh$.
}
\end{ex}

% ===== Câu 84 =====
% ID: L12C2B8-104-TN
% Mức: VDC
\begin{ex}
% ID: L12C2B8-104-TN
% Mức: VDC
Trong dạng “Khối khí chịu áp suất của cột chất lỏng”, Trong ống có cột chất lỏng, áp suất khí bị giam được xác định bằng
\choice
{\True cân bằng áp suất tại cùng độ cao}
{chỉ áp suất khí quyển}
{chỉ $\rho gh$}
{luôn bằng không}
\loigiai{
Đáp án A. Dùng $p_{khí}=p_0\pm\rho gh$ tùy cách bố trí.
}
\end{ex}

% ===== Câu 106 =====
% ID: L12C2B8-105-DS
% Mức: VD
\begin{ex}
% ID: L12C2B8-105-DS
% Mức: VD
Trong dạng “Đồ thị đẳng nhiệt trên các hệ tọa độ”, Nêu cách xác định áp suất cột khí bị giam bởi chất lỏng.
\choiceTF

\loigiai{
Lập cân bằng áp suất tại cùng một độ cao trong chất lỏng; tùy vị trí mặt thoáng, $p_{khí}=p_0\pm\rho gh$.
}
\end{ex}

% ===== Câu 124 =====
% ID: L12C2B8-106-TN
% Mức: VDC
\begin{ex}
% ID: L12C2B8-106-TN
% Mức: VDC
Trong dạng “Nhận biết quá trình đẳng nhiệt và định luật Boyle”, Áp suất tuyệt đối của cột khí bị giam bằng
\choice
{\True áp suất khí quyển cộng hoặc trừ áp suất cột chất lỏng tùy bố trí}
{luôn bằng áp suất khí quyển}
{chỉ bằng áp suất cột chất lỏng}
{bằng không}
\loigiai{
Đáp án A. Cân bằng áp suất phải xét cả khí quyển và chênh áp thủy tĩnh.
}
\end{ex}

% ===== Câu 144 =====
% ID: L12C2B8-107-TN
% Mức: VDC
\begin{ex}
% ID: L12C2B8-107-TN
% Mức: VDC
Trong dạng “Tính áp suất, thể tích trong quá trình đẳng nhiệt”, Áp suất tuyệt đối của cột khí bị giam bằng
\choice
{\True áp suất khí quyển cộng hoặc trừ áp suất cột chất lỏng tùy bố trí}
{luôn bằng áp suất khí quyển}
{chỉ bằng áp suất cột chất lỏng}
{bằng không}
\loigiai{
Đáp án A. Cân bằng áp suất phải xét cả khí quyển và chênh áp thủy tĩnh.
}
\end{ex}

% ===== Câu 154 =====
% ID: L12C2B8-108-TN
% Mức: VDC
\begin{ex}
% ID: L12C2B8-108-TN
% Mức: VDC
Trong dạng “Đồ thị đẳng nhiệt trên các hệ tọa độ”, Áp suất tuyệt đối của cột khí bị giam bằng
\choice
{\True áp suất khí quyển cộng hoặc trừ áp suất cột chất lỏng tùy bố trí}
{luôn bằng áp suất khí quyển}
{chỉ bằng áp suất cột chất lỏng}
{bằng không}
\loigiai{
Đáp án A. Cân bằng áp suất phải xét cả khí quyển và chênh áp thủy tĩnh.
}
\end{ex}

% ===== Câu 160 =====
% ID: L12C2B8-109-TN
% Mức: VDC
\begin{ex}
% ID: L12C2B8-109-TN
% Mức: VDC
Trong ống có cột chất lỏng, áp suất khí bị giam được xác định bằng
\choice
{\True cân bằng áp suất tại cùng độ cao}
{chỉ áp suất khí quyển}
{chỉ $\rho gh$}
{luôn bằng không}
\loigiai{
Đáp án A. Dùng $p_{khí}=p_0\pm\rho gh$ tùy cách bố trí.
}
\end{ex}
